\documentclass[11pt]{article}
\usepackage{hwtools}

\newcommand{\C}{\mathbb{C}}
\newcommand{\R}{\mathbb{R}}
\newcommand{\Q}{\mathbb{Q}}
\newcommand{\Z}{\mathbb{Z}}
\newcommand{\F}{\mathbb{F}}

\title{Homework 2}
\author{Asher Lucas-Cuddeback}

\begin{document}
\maketitle
\setcounter{problemcounter}{8}

\begin{problem}[Orbits, Stabilizers, and Normality]
  Consider the group action of the dihedral group
  $D_n \curvearrowright \{n \text{ vertices of a regular } n\text{-gon}\}$.
  For example, $D_3$ acts on the 3 vertices of the equilateral triangle,
  $D_4$ acts on the 4 vertices of the square, etc.
  \begin{parts}

  \item Describe the space of orbits of this action.
    \begin{solution}

    \end{solution}

  \item Describe the orbit as $D_n/H$ for some subgroup $H$. Is $H \le D_n$
    normal?
    \begin{solution}

    \end{solution}

  \item Show that the cyclic subgroup $C_n \le D_n$ is normal, and thus
    $D_n/C_n$ actually forms a group. What group is it isomorphic to? What
    space does $D_n/C_n$ naturally act on?
    \begin{solution}

    \end{solution}

  \end{parts}
\end{problem}
\newpage

\begin{problem}[Geometry Emerges from $G/H$ (Finite Examples)]
  Explain the following equivalences of sets.
  \begin{parts}

  \item $\{n \text{ ordered points}\} \cong S_n$
    \begin{solution}

    \end{solution}

  \item $\{n \text{ unordered points}\} \cong S_n/S_{n-1}$
    \begin{solution}

    \end{solution}

  \item $\{\text{ordered pairs of distinct elements from } \{1,\dots,n\}\}
    \cong S_n/S_{n-2}$
    \begin{solution}

    \end{solution}

  \item $\{\text{unordered pairs of distinct elements from } \{1,\dots,n\}\}
    \cong S_n/(S_{n-2}\times S_2)$
    \begin{solution}

    \end{solution}

  \end{parts}
\end{problem}
\newpage

\begin{problem}[Geometry Emerges from $G/H$ (Continuous Examples)]
  The four famous, and historically important geometries are: Euclidean,
  Spherical, Hyperbolic, and Projective. Felix Klein's Erlangen Program was
  the idea that the study of geometry is inextricably linked with the study of
  symmetry groups. The relevant symmetry groups required to produce these
  geometries are the Euclidean group, the Orthogonal group, and real General
  Linear group.

  Explain the following equivalences of two-dimensional geometries.
  \begin{parts}

  \item $\mathbb{E}^2 = \{\text{the plane}\} \cong E(2)/O(2)$
    \begin{solution}

    \end{solution}

  \item $\mathbb{S}^2 = \{\text{the 2-sphere}\} \cong O(3)/O(2) \cong
    SO(3)/SO(2)$
    \begin{solution}

    \end{solution}

  \item $\mathbb{H}^2 = \{\text{the upper half plane}\} \cong
    SL_2(\R)/SO(2) \cong PSL_2(\R)/PSO(2)$
    \begin{solution}

    \end{solution}

  \item $\mathbb{P}^2 = \{\text{lines through the origin in } \R^3\}
    \cong GL_3(\R)/P$

    (describe the subgroup $P \le GL_2(\R)$, it is called a parabolic
    subgroup)
    \begin{solution}

    \end{solution}

  \end{parts}
\end{problem}

\end{document}
